\section{Arqueología de algoritmos: cómo los artículos antiguos vuelven a brillar con el hardware nuevo}

Yang Songlin (杨松琳) enfatiza la «arqueología»: leer artículos antiguos, entender mecanismos antiguos y después recombinarlos frente a problemas nuevos y hardware nuevo. Muchos mecanismos ya habían aparecido en 2016, 2020 y 2021, pero en su momento faltaban el hardware adecuado, algoritmos paralelos, implementaciones de código abierto o validación a gran escala, así que quedaron sepultados en el mar de la literatura. DataNet, Delta Rule, Gated Linear Attention y otros pertenecen a este tipo de pistas reactivadas.

\begin{figure}[H]
\centering
\includegraphics[width=0.96\textwidth]{algorithm-archaeology.png}
\caption{Ruta de la arqueología de algoritmos: artículos antiguos, mecanismos antiguos y hardware nuevo se recombinan. Diagrama conceptual propio, elaborado a partir de la conversación de 01:23:12--01:29:50.}
\end{figure}

\begin{knowledgebox}{Lectura de la figura: la arqueología no es nostalgia, sino la búsqueda de herramientas reutilizables}
Entre Transformer, Sparse, DeltaNet, Gated Linear y KDA no hay un progreso en línea recta, sino mecanismos antiguos que se reevalúan en un contexto nuevo. Muchas ideas no aparecen por primera vez, sino que por primera vez reúnen las condiciones para hacer scale up.
\end{knowledgebox}

\subsection{Por qué los trabajos antiguos quedan sepultados}

Esta sección explica por qué la «arqueología» no es un adorno. Muchos artículos antiguos no estaban equivocados; en su momento carecían de las condiciones de ingeniería que les permitieran convertirse en la corriente principal.

Hay tres posibles razones por las que un trabajo antiguo queda sepultado. Primera: los algoritmos complementarios y los kernels no están maduros, de modo que a otros les resulta difícil reproducirlo o escalarlo. Segunda: el hardware y las tareas de la época no lo necesitaban, así que su valor no se hizo visible. Tercera: el código abierto es difícil de usar y la comunidad no puede darle continuidad. Yang Songlin enfatiza hacer bien el código para que la técnica pueda transmitirse; en realidad, esto es un componente importante del impacto de la investigación en arquitecturas.

\begin{warningbox}{Que funcione a pequeña escala no significa que funcione a gran escala}
Existen muchas variantes de arquitectura, y muchos algoritmos funcionan en tareas a pequeña escala, pero fallan en el preentrenamiento a gran escala. El Scaling Ladder, la implementación en hardware y la reproducción en código abierto son mecanismos importantes para evitar la ilusión de la pequeña escala.
\end{warningbox}

\subsection{Criterio de investigación: saber qué vale la pena hacer}

Yang Songlin dice que el verdadero gran reto no es el problema técnico en sí, sino no saber qué hacer. Su método consiste en leer primero todos los artículos del campo que vale la pena revisar, formar un mapa de problemas y luego juzgar qué mecanismo antiguo tiene valor hoy. Por ejemplo, DataNet ya existía en 2021, pero no tenía garantía de eficiencia en hardware; cuando las necesidades de las tareas y los algoritmos paralelos maduran, puede volver a ser un componente utilizable.

\begin{importantbox}{El criterio en la investigación de algoritmos}
La buena investigación en arquitecturas no consiste en combinar módulos al azar, sino en saber qué herramientas existieron en la historia, por qué no tuvieron éxito, qué restricciones han cambiado hoy y cómo usar algoritmos y hardware nuevos para hacer scale up de las ideas antiguas.
\end{importantbox}

\subsection{La implementación de código abierto también es una contribución de investigación}

Flash Linear Attention es importante no solo por el artículo, sino porque permite que investigadores y la industria prueben con facilidad la atención lineal. Si una idea no tiene una implementación utilizable, a la investigación posterior le cuesta darle seguimiento; si la implementación es fácil de usar, la línea técnica se transmite y se valida con más facilidad.

\begin{knowledgebox}{Tres pasos del artículo a la línea de investigación}
El primer paso es la idea: que tenga sentido en lo matemático o en el mecanismo. El segundo paso es la implementación: si puede escribirse como código utilizable, reproducible y optimizable. El tercer paso es el scaling: si se sostiene con modelos más grandes, tareas más realistas y condiciones más favorables para el hardware.
\end{knowledgebox}

\subsection{Resumen de la sección}

La arqueología de algoritmos ofrece un método de investigación: no hay que perseguir solo la buzzword más reciente, sino leer artículos antiguos, entender los mecanismos y juzgar qué ideas pueden volver a brillar con el hardware y las tareas de hoy.
